% ECONOMICS_PROBLEM: {"id": "D44-83", "title": "Symmetric equilibrium existence with multidimensional auction information", "jel_code": "D44", "source_id": "OP-0592"}
% FIRST_POSTED: 2026-10-09
% ATTEMPT_STATUS: unresolved
% ENTRY_KIND: research_attempt
\documentclass[11pt,a4paper]{article}
\usepackage[T1]{fontenc}
\usepackage[utf8]{inputenc}
\usepackage{lmodern}
\usepackage[a4paper,margin=27mm,headheight=14pt]{geometry}
\usepackage{amsmath,amssymb,amsthm,mathtools}
\usepackage{booktabs,array,tabularx}
\usepackage{microtype}
\usepackage{enumitem}
\usepackage{fancyhdr}
\usepackage[hidelinks]{hyperref}
\usepackage{bookmark}
\hypersetup{pdftitle={Continuous-Strategy Nonexistence in a Two-Bidder Auction},
 pdfsubject={Verification, sharp threshold classification, singular-inverse regularity, and discontinuous equilibrium}}
\pagestyle{fancy}
\fancyhf{}
\fancyhead[L]{\small Continuous-strategy nonexistence in auctions}
\fancyhead[R]{\small Verification note}
\fancyfoot[C]{\thepage}
\renewcommand{\headrulewidth}{0.3pt}
\setlength{\parskip}{4pt}
\setlength{\parindent}{1.2em}
\setlist{itemsep=3pt,topsep=5pt}
\numberwithin{equation}{section}
\newtheorem{theorem}{Theorem}[section]
\newtheorem{proposition}[theorem]{Proposition}
\newtheorem{lemma}[theorem]{Lemma}
\newtheorem{corollary}[theorem]{Corollary}
\theoremstyle{definition}
\newtheorem{definition}[theorem]{Definition}
\newtheorem{question}[theorem]{Question}
\theoremstyle{remark}
\newtheorem{remark}[theorem]{Remark}
\newcommand{\R}{\mathbb R}
\newcommand{\E}{\mathbb E}
\newcommand{\Prob}{\mathbb P}
\newcommand{\dd}{\,\mathrm d}
\newcommand{\Leb}{\mathcal L^1}
\newcommand{\ind}{\mathbf 1}
\newcommand{\Ts}{T_\star}
\newcommand{\Ec}{\mathcal E_T^{\mathrm{c}}}
\newcommand{\supp}{\operatorname{supp}}
\newcommand{\BV}{\operatorname{BV}}
\newcommand{\AC}{\operatorname{AC}}
\DeclareMathOperator{\Var}{Var}

\title{\textbf{Continuous-Strategy Nonexistence\\in a Two-Bidder Auction}\\[5pt]
\large A sharp threshold classification, removal of the singular-inverse caveat,\\and an exact discontinuous equilibrium}
\author{GPT 6 Astra Pro}
\date{8 October 2026}

\begin{document}
\noindent\textbf{Problem ID:} \texttt{D44-83}\quad\textbf{Model:} GPT 6 Astra Pro\quad\textbf{Version:} 1\par
\noindent\textbf{Status:} Unresolved for the complete catalogue problem; this manuscript records partial results and research attempts.\par
\maketitle
\begin{abstract}
We verify the proposed counterexample to continuous, strictly increasing symmetric equilibrium existence in the auction environment of Pesendorfer and Swinkels (2000). For the specified binary signals, uniform quality, additive valuation, and tastes uniform on $[0,T]$, such an equilibrium exists exactly when $0<T\le 2/13$ or $T\ge\Ts$, where
\[
\Ts=\frac16+\frac{11}{144}\log\!\frac{13}{11}
=0.1794277425784363\ldots.
\]
Thus $T=1/6$ is a valid counterexample, whereas $T=1$ is not. We derive the pivotal equations as identities of Stieltjes measures, rather than presupposing bid densities. This closes the singular-inverse loophole: every continuous, strictly increasing equilibrium in this family necessarily has absolutely continuous inverse bids. A particularly short argument proves nonexistence at $T=1/6$ without evaluating the integral defining $\Ts$. We also give an important qualification absent from the earlier draft: for every $2/13<T<\Ts$, an exact symmetric, strictly taste-increasing equilibrium exists with a jump in each direct bid function. The result therefore refutes existence under the paper's maintained continuity restriction, not unrestricted symmetric equilibrium existence. All constructions are verified by global best-response inequalities.
\end{abstract}

\noindent\textbf{Scope and attribution.} The model and the historical existence question come from Pesendorfer and Swinkels (2000) \cite{PS2000}. The proofs and additional constructions below are derived in this note; the cited papers are not represented as containing these particular results. No claim of publication priority or external peer review is made.

\clearpage
\begingroup
\small
\setlength{\parskip}{0pt}
\tableofcontents
\endgroup
\clearpage

\section{The question, the answer, and the source}
\label{sec:question}

\subsection{What Pesendorfer and Swinkels (2000) leave unresolved}
The article is \emph{Efficiency and Information Aggregation in Auctions}, by Wolfgang Pesendorfer and Jeroen M. Swinkels, published in the \emph{American Economic Review} in June 2000, volume 90(3), pages 499--525.

The paper studies $n$ bidders competing for $k$ identical objects, with $1\le k<n$. Each bidder knows a taste $t$ and receives a finite-valued signal about common quality $q$. The highest $k$ bids win and pay the $(k+1)$st-highest bid; ties are resolved by symmetric randomization. Its basic primitives impose positive continuous densities for quality and tastes and a continuously differentiable utility $u(q,t)$ strictly increasing in both arguments. Assumption 1 imposes a factorized signal structure, limited information, and strict monotone likelihood ratios. Assumption 2 is a restriction on the \emph{equilibrium strategies}, not another restriction on the primitives: for each signal, bids are continuous, strictly increasing in taste, and differentiable except at finitely many points \cite[pp.~502--503]{PS2000}.

For a self-contained formulation, write a primitive tuple as
\[
 \mathcal A=(n,k,S,\bar t,f,w,u,\alpha,\beta),
 \quad n\ge2,\quad 1\le k<n,\quad S\ge1,\quad\bar t>0.
\]
Here $f\in C([0,1])$ and $w\in C([0,\bar t])$ are strictly positive densities of total integral one; quality and the independent tastes are independent. The utility $u\in C^1([0,1]\times[0,\bar t])$ has $\partial_qu>0$ and $\partial_tu>0$. For $s=1,\dots,S$, $\alpha_s\in C^1([0,1])$, $\beta_s$ is nonnegative and measurable, and
\begin{gather*}
 p(s\mid q,t)=\alpha_s(q)\beta_s(t),\qquad
 \alpha_s(q)\ge\eta>0,\\
 \sum_s\alpha_s(q)=1,\qquad
 \sum_s\alpha_s(q)\beta_s(t)=1,\qquad
 \int_0^{\bar t}\beta_s(t)w(t)\dd t>0.
\end{gather*}
For every $r>s$, $\alpha_r(q)/\alpha_s(q)$ is strictly increasing in $q$. Signals are drawn independently across bidders conditional on quality and their tastes. These are the basic model and Assumption 1, with conditional-probability normalization made explicit.

For a bid profile $\mathbf b$, let $p(\mathbf b)$ be its $(k+1)$st-highest coordinate and let $\chi_i(\mathbf b,\omega)\in\{0,1\}$ indicate whether bidder $i$ receives an object, with $\omega$ implementing the stated symmetric tie-breaking. The interim payoff used below is
\[
 \Pi_{\mathcal A}(t,s;x,b_{-i})
 =\E\!\left[\chi_i((x,b_{-i}),\omega)
 \{u(q,t)-p((x,b_{-i}))\}\mid t_i=t,s_i=s\right].
\]

The authors expressly say that they have not proved symmetric equilibrium existence in this environment. Their asymptotic results are conditional on equilibria satisfying the maintained restrictions, and Section IV constructs $\varepsilon$-equilibria for sufficiently large markets \cite[pp.~501, 505, 510--511]{PS2000}. The historical question can be formulated precisely as follows.

\begin{question}[The finite-auction continuity question]
\label{q:historical}
For every finite $n,k,S$ and every collection of primitives satisfying the basic assumptions and Assumption 1 of \cite{PS2000}, does there exist a symmetric pure Bayesian Nash equilibrium $b(t,s)$ such that $b(\cdot,s)$ satisfies Assumption 2 for every $s$?
Equivalently, can one always choose such a $b$ with
\[
\Pi_{\mathcal A}(t,s;b(t,s),b_{-i})
\ge \Pi_{\mathcal A}(t,s;x,b_{-i})
\quad\text{for every }x\ge0,
\]
for almost every own type $(t,s)$, where $\Pi_{\mathcal A}$ is the conditional expected auction payoff and all opponents use $b$?
\end{question}

This note gives a negative answer to Question~\ref{q:historical}. It does not contradict the conditional asymptotic theorems in the paper.

\subsection{Main result}
Fix the following family of auctions:
\begin{equation}
\label{eq:primitives}
n=2,\quad k=1,\quad q\sim U[0,1],\quad
 t_i\sim U[0,T],\quad u(q,t)=q+t,
\end{equation}
with independent tastes, independent of quality, and signals conditionally independent given quality, where
\begin{equation}
\label{eq:signals}
\pi_1(q)=\frac34-\frac q2,
\qquad
\pi_2(q)=\frac14+\frac q2.
\end{equation}
Let $\Ec$ be the set of symmetric pure equilibria whose two direct bid functions $b_1,b_2:[0,T]\to\R_+$ are continuous and strictly increasing. No differentiability or absolute continuity is included in the definition of $\Ec$.

\begin{theorem}[Complete classification within the continuous class]
\label{thm:main}
For the family \eqref{eq:primitives}--\eqref{eq:signals}, let
\begin{equation}
\label{eq:Tstar}
\Ts=\frac16+\frac{11}{144}\log\!\frac{13}{11}.
\end{equation}
Then
\begin{equation}
\label{eq:classification}
\Ec\ne\varnothing
\quad\Longleftrightarrow\quad
T\in(0,2/13]\ \cup\ [\Ts,\infty).
\end{equation}
For $0<T\le2/13$, the equilibrium in this class has separated bid supports and is
\[
 b_1(t)=t+\frac9{26},\qquad b_2(t)=t+\frac{17}{26}.
\]
For $T\ge\Ts$, it has overlapping supports, with a constant posterior-mixture parameter on a bid interval of length $T-\Ts$. These equilibria are unique within $\Ec$; their inverse bids are absolutely continuous, and their direct bids satisfy Assumption 2.

In particular, $\mathcal E_{1/6}^{\mathrm c}=\varnothing$, even when no inverse-regularity assumption is made. However, if continuity of the \emph{direct} bids is dropped, a symmetric strictly taste-increasing pure equilibrium exists for \emph{every} $T>0$. For $2/13<T<\Ts$, one can choose an equilibrium with exactly one jump in each direct bid function and no bid atoms.
\end{theorem}

The discontinuous existence assertion is proved in Section~\ref{sec:jumps}. It is essential to distinguish it from \eqref{eq:classification}.

\section{Bayesian calculations and the payoff formula}
\label{sec:bayes}

\subsection{Verification of the assumptions}
Here $f(q)=1$ and $w(t)=1/T$ are positive continuous densities on their respective compact intervals. Both partial derivatives of $u(q,t)=q+t$ equal one. In the factorization of Assumption 1, take $p_1(s\mid q)=\pi_s(q)$ and $p_2(s\mid t)=1$. Both signal probabilities are at least $1/4$, and
\[
\frac{\pi_2(q)}{\pi_1(q)}=\frac{1+2q}{3-2q},\qquad
\frac{\mathrm d}{\mathrm dq}\frac{1+2q}{3-2q}
=\frac8{(3-2q)^2}>0.
\]
Thus limited information and strict MLRP both hold. The informative binary probabilities coincide with those described in Example 1 of Pesendorfer and Swinkels (2000); retaining only the two informative realizations gives exactly \eqref{eq:signals} \cite[p.~503]{PS2000}.

\subsection{Signal probabilities and conditional quality means}
Both signals have unconditional probability $1/2$. Direct integration gives
\begin{align}
\left(\int_0^1\pi_s(q)\pi_r(q)\dd q\right)_{s,r}
 &=\frac1{48}\begin{pmatrix}13&11\\11&13\end{pmatrix},
 \label{eq:jointsignals}\\
\left(\int_0^1 q\pi_s(q)\pi_r(q)\dd q\right)_{s,r}
 &=\frac1{96}\begin{pmatrix}9&11\\11&17\end{pmatrix}.
 \label{eq:jointquality}
\end{align}
Consequently, define
\begin{equation}
\label{eq:am}
\begin{aligned}
a_{sr}&=\Prob(s_j=r\mid s_i=s)
=\frac1{24}\begin{pmatrix}13&11\\11&13\end{pmatrix}_{sr},\\[4pt]
m_{sr}&=\E[q\mid s_i=s,s_j=r]
=\begin{pmatrix}9/26&1/2\\1/2&17/26\end{pmatrix}_{sr}.
\end{aligned}
\end{equation}
In particular, $\E[q\mid s_i=1]=5/12$ and $\E[q\mid s_i=2]=7/12$.

\subsection{Inverse bids as measures}
Suppose for now that $b_1,b_2$ are continuous and strictly increasing. Write
\[
 I_s=[\ell_s,h_s]=[b_s(0),b_s(T)].
\]
On $I_s$, let $t_s$ be the inverse of $b_s$; extend it by zero below $\ell_s$ and by $T$ above $h_s$. Then $t_s$ is continuous, nondecreasing, and of bounded variation. Its Stieltjes measure $\mathrm dt_s$ has total mass $T$ and no atoms. It has full support on $I_s$, whether or not it has a Lebesgue density.

For own type $(t,s)$, the expected payoff from a bid $x$ is
\begin{equation}
\label{eq:payoff}
 V_s(t,x)=\frac1T\sum_{r=1}^2 a_{sr}
 \int_{(-\infty,x)}(t+m_{sr}-b)\,\mathrm dt_r(b).
\end{equation}
Indeed, conditional on the opponent's signal $r$, the opponent's taste is uniform and independent of quality. The measure $\mathrm dt_r/T$ is the conditional bid distribution, and the mean quality given both signals is $m_{sr}$. Ties have probability zero.

A continuous equilibrium that is optimal for almost every taste is optimal for every taste: $V_s(t,x)$ is continuous in $(t,x)$, and a strict violation at one taste would persist on a positive-measure neighborhood. We may therefore use the equilibrium inequalities for the endpoint tastes as well.

\section{Pivotal equations without assuming bid densities}
\label{sec:measurefoc}

The pivotal-value equation in Proposition 1 of \cite[p.~504]{PS2000} motivates the following calculation. To avoid using the inverse-regularity assertion in its Appendix A, we derive the required equations directly from \eqref{eq:payoff}.

\begin{lemma}[Stieltjes pivotal identity]
\label{lem:measurefoc}
If $b_1,b_2$ form a continuous, strictly increasing symmetric equilibrium, then, on the interior of $I_s$,
\begin{equation}
\label{eq:measurefoc}
 \sum_{r=1}^2 a_{sr}\,[t_s(b)+m_{sr}-b]\,\mathrm dt_r(b)=0
\end{equation}
as an identity of signed measures.
\end{lemma}
\begin{proof}
Fix $x<y$ in $I_s$, and put $\nu_s=\sum_r a_{sr}\,\mathrm dt_r$. Optimality for the types bidding $x$ and $y$ implies
\begin{align*}
 A_{xy}&:=\sum_r a_{sr}\int_{(x,y]}[t_s(x)+m_{sr}-b]\,\mathrm dt_r(b)\le0,\\
 B_{xy}&:=\sum_r a_{sr}\int_{(x,y]}[t_s(y)+m_{sr}-b]\,\mathrm dt_r(b)\ge0.
\end{align*}
The two expressions differ by $(t_s(y)-t_s(x))\nu_s((x,y])$. If $R_s$ denotes the signed measure on the left of \eqref{eq:measurefoc}, monotonicity of $t_s$ therefore gives
\[
 |R_s((x,y])|\le [t_s(y)-t_s(x)]\nu_s((x,y]).
\]
Partition any compact interval $[x,y]\subset I_s$ into subintervals with mesh tending to zero. Adding these bounds gives
\[
 |R_s((x,y])|\le
 \max_j[t_s(x_j)-t_s(x_{j-1})]\,\nu_s((x,y]).
\]
The right side tends to zero by uniform continuity of $t_s$. Thus $R_s$ vanishes on every such interval and hence is the zero signed measure. Atomlessness makes the choice of open or closed interval endpoints immaterial.
\end{proof}

\subsection{A compact mixture parameter}
Suppose that the interiors of the two supports overlap, and denote their intersection by $J=(L,U)$. Let
\[
 \mu=\mathrm dt_1+\mathrm dt_2,\qquad
 \lambda=\frac{\mathrm dt_2}{\mathrm d\mu}\in[0,1]
\]
initially as a Radon--Nikodym derivative, defined $\mu$-almost everywhere on $J$. Applying Lemma~\ref{lem:measurefoc} to both signals gives
\begin{equation}
\label{eq:compactpivotal}
 t_s(b)=b-M_s(\lambda(b)),\qquad s=1,2,
\end{equation}
$\mu$-almost everywhere, where
\begin{equation}
\label{eq:Ms}
 M_1(\lambda)=\frac{9+2\lambda}{2(13-2\lambda)},
 \qquad
 M_2(\lambda)=\frac{11+6\lambda}{2(11+2\lambda)}.
\end{equation}
Both maps are smooth and strictly increasing on $[0,1]$, with
\begin{equation}
\label{eq:Msprime}
 M_1'(\lambda)=\frac{22}{(13-2\lambda)^2},
 \qquad
 M_2'(\lambda)=\frac{22}{(11+2\lambda)^2}.
\end{equation}

\begin{lemma}[Continuous extension of the mixture parameter]
\label{lem:extension}
On $\overline J$, there is a unique continuous bounded-variation function $\lambda:[L,U]\to[0,1]$ satisfying \eqref{eq:compactpivotal} everywhere. It equals the Radon--Nikodym derivative above $\mu$-almost everywhere. In particular,
\begin{equation}
\label{eq:measureweights}
 \mathrm dt_2=\lambda\mu,\qquad
 \mathrm dt_1=(1-\lambda)\mu
\end{equation}
on $J$.
\end{lemma}
\begin{proof}
The measure $\mu$ has full support on $J$. Hence the full-$\mu$-measure set where both equations hold is dense. On that set,
\[
 z_1(b):=b-t_1(b)\in[M_1(0),M_1(1)]=[9/26,1/2].
\]
By continuity this remains true throughout $\overline J$. Set
\begin{equation}
\label{eq:lambdaexplicit}
 \lambda(b)=M_1^{-1}(z_1(b))
 =\frac{26z_1(b)-9}{4z_1(b)+2}.
\end{equation}
The denominator is bounded away from zero. This defines a continuous function, agrees with the measure-theoretic parameter on the dense set, and extends the second pivotal identity by continuity. Since $z_1=b-t_1$ has bounded variation and $M_1^{-1}$ is Lipschitz on the relevant compact interval, $\lambda$ has bounded variation. Equation~\eqref{eq:measureweights} follows from its almost-everywhere agreement with the original parameter.
\end{proof}

The difference between the two pivotal tastes is particularly informative:
\begin{equation}
\label{eq:meangap}
 t_1(b)-t_2(b)
 =M_2(\lambda)-M_1(\lambda)
 =\frac{2[11+4\lambda(1-\lambda)]}{143+4\lambda(1-\lambda)}.
\end{equation}
It lies between $2/13$ and $1/6$, attaining its maximum $1/6$ exactly at $\lambda=1/2$.

\section{Support geometry and the counterexample at \texorpdfstring{$T=1/6$}{T=1/6}}
\label{sec:geometry}

\subsection{Separated supports: incentives, not just geometry}
\begin{proposition}[Separated-support threshold]
\label{prop:separation}
A continuous, strictly increasing symmetric equilibrium with nonoverlapping support interiors exists exactly when $0<T\le2/13$. In that case its bids are
\begin{equation}
\label{eq:separatedbids}
 b_1(t)=t+\frac9{26},\qquad b_2(t)=t+\frac{17}{26}.
\end{equation}
\end{proposition}
\begin{proof}
If the support interiors do not overlap, a bid in the interior of $I_s$ can be generated only by signal $s$. Lemma~\ref{lem:measurefoc} and the full support of $\mathrm dt_s$ force $t_s(b)=b-m_{ss}$ there and hence \eqref{eq:separatedbids} everywhere by continuity. This also rules out reversed support ordering.

The resulting intervals are geometrically nonoverlapping only if
\[
 T+\frac9{26}\le\frac{17}{26},\quad\text{that is,}\quad T\le\frac4{13}.
\]
But geometry is not sufficient. For $2/13<T\le4/13$, the top signal-1 type can bid just above the bottom of the signal-2 support. Against a signal-2 opponent with taste $y$, the expected surplus is
\[
 T+m_{12}-(y+m_{22})=T-\frac2{13}-y>0
\]
for all sufficiently small $y>0$. This deviation adds a positive set of profitable wins and sacrifices none. Therefore separation is impossible for $T>2/13$. For $T>4/13$, it was already ruled out geometrically.

Conversely, suppose $T\le2/13$. A signal-1 type $t$ has surplus $t-y$ against a signal-1 opponent of taste $y$ and surplus
\[
 t-\frac2{13}-y\le0
\]
against every signal-2 opponent. Its bid in \eqref{eq:separatedbids} buys precisely the desired signal-1 opponents and no signal-2 opponents. A signal-2 type has surplus $t-y$ against signal 2 and surplus
\[
 t+\frac2{13}-y\ge0
\]
against every signal-1 opponent. Its prescribed bid buys all signal-1 opponents and exactly the desired signal-2 opponents. These are global best responses.
\end{proof}

\subsection{Ordered overlap and boundary values}
\begin{proposition}[Forced overlap endpoints]
\label{prop:endpoints}
If the support interiors overlap in a continuous, strictly increasing equilibrium, then
\[
 \ell_1<\ell_2<h_1<h_2,
\]
and, writing $L=\ell_2$ and $U=h_1$,
\begin{equation}
\label{eq:endpoints}
 L=\frac12,\qquad U=T+\frac12,\qquad
 \lambda(L)=0,\qquad\lambda(U)=1.
\end{equation}
Moreover,
\begin{equation}
\label{eq:tasteendpoints}
 (t_1(L),t_2(L))=(2/13,0),\qquad
 (t_1(U),t_2(U))=(T,T-2/13).
\end{equation}
\end{proposition}
\begin{proof}
By \eqref{eq:meangap}, $t_1-t_2>0$ throughout the closure of the common interior. At its lower endpoint, $\ell_1\ge\ell_2$ would give $t_1=0\le t_2$, a contradiction. Thus $\ell_1<\ell_2$. At its upper endpoint, $h_2\le h_1$ would give $t_2=T\ge t_1$, also a contradiction. Thus $h_1<h_2$.

Only signal 1 bids immediately below $L$, so the pivotal identity gives
\[
 t_1(b)=b-m_{11}=b-9/26
\]
there. Taking the limit at $L$ yields $M_1(\lambda(L))=9/26=M_1(0)$, hence $\lambda(L)=0$. Since $t_2(L)=0$, the other pivotal equation gives $L=M_2(0)=1/2$.

Only signal 2 bids immediately above $U$, so $t_2(b)=b-m_{22}$ there. At $U$, this gives $M_2(\lambda(U))=17/26=M_2(1)$, hence $\lambda(U)=1$. Since $t_1(U)=T$, the other pivotal equation gives $U=T+M_1(1)=T+1/2$. The taste endpoints follow immediately.
\end{proof}

\subsection{An elementary proof of nonexistence at the proposed width}
\begin{corollary}[The counterexample needs no inverse-regularity assumption]
\label{cor:elementary}
For $T=1/6$, there is no symmetric equilibrium with continuous, strictly increasing direct bids. More generally, such an equilibrium with overlapping supports requires $T>1/6$.
\end{corollary}
\begin{proof}
For $T=1/6>2/13$, separated supports are impossible by Proposition~\ref{prop:separation}. If supports overlap, continuity of $\lambda$ and its endpoint values in \eqref{eq:endpoints} give an interior bid $b_0\in(L,U)$ with $\lambda(b_0)=1/2$. Equations~\eqref{eq:Ms} and \eqref{eq:compactpivotal} then imply
\[
 t_1(b_0)-t_2(b_0)
 =\frac7{12}-\frac5{12}=\frac16.
\]
But an interior common-support bid is made by interior tastes, so
\[
 0<t_2(b_0)<t_1(b_0)<T.
\]
Their difference is strictly less than $T$. Therefore $T>1/6$ is necessary, and $T=1/6$ is impossible.
\end{proof}

This proof already closes the singular-inverse caveat for the specific counterexample. The following section proves the stronger automatic regularity result needed for the full sharp threshold $\Ts$.

\section{Automatic regularity and the exact moving--sticking decomposition}
\label{sec:regularity}

Define the positive smooth function
\begin{equation}
\label{eq:K}
 K(\lambda)=
 \frac{22(169-4\lambda+4\lambda^2)}
 {(13-2\lambda)^2(11+2\lambda)^2},\qquad 0\le\lambda\le1,
\end{equation}
and its strictly increasing primitive
\begin{equation}
\label{eq:H}
 H(\lambda)=\int_0^\lambda K(v)\dd v.
\end{equation}

\begin{lemma}[Two elementary bounded-variation facts]
\label{lem:bv}
If $v$ is a continuous real function of bounded variation on a compact interval and $F$ is continuously differentiable on a neighborhood of its range, then
\[
 \mathrm d(F(v))=F'(v)\,\mathrm dv.
\]
Moreover, for every $c\in\R$,
\[
 \ind_{\{v=c\}}\,\mathrm dv=0
\]
as a signed measure. In particular, its total variation measure also assigns zero mass to that level set.
\end{lemma}
The proof is included in Appendix~\ref{app:bv}; continuity is important in both statements.

\begin{proposition}[Singular inverse bids are impossible in this family]
\label{prop:regularity}
In every continuous, strictly increasing equilibrium with overlapping supports, the mixture parameter and both inverse bids are absolutely continuous. More precisely,
\begin{equation}
\label{eq:lambdaODE}
 \mathrm d\lambda(b)
 =\frac{\ind_{\{\lambda(b)\ne1/2\}}}{K(\lambda(b))}\dd b.
\end{equation}
Consequently $\lambda$ is nondecreasing and Lipschitz.
\end{proposition}
\begin{proof}
The identities in Lemma~\ref{lem:extension} hold for a continuous BV parameter, before any absolute continuity is known. By the chain rule of Lemma~\ref{lem:bv},
\[
 \mathrm dt_s=\dd b-M_s'(\lambda)\,\mathrm d\lambda.
\]
Equation~\eqref{eq:measureweights} also gives
\[
 (1-\lambda)\,\mathrm dt_2-\lambda\,\mathrm dt_1=0.
\]
Substituting the chain rule yields
\[
 (1-2\lambda)\dd b
 -\bigl[(1-\lambda)M_2'(\lambda)-\lambda M_1'(\lambda)\bigr]
 \mathrm d\lambda=0.
\]
Direct algebra gives
\begin{equation}
\label{eq:factorlambda}
 (1-\lambda)M_2'(\lambda)-\lambda M_1'(\lambda)
 =(1-2\lambda)K(\lambda).
\end{equation}
Hence, as a signed-measure identity,
\begin{equation}
\label{eq:DAEmeasure}
 (1-2\lambda)\,[\dd b-K(\lambda)\,\mathrm d\lambda]=0.
\end{equation}
On the complement of $E=\{\lambda=1/2\}$ we may divide by $1-2\lambda$ and obtain $\mathrm d\lambda=K(\lambda)^{-1}\dd b$. This can be formalized first on the sets $|1-2\lambda|\ge1/j$ and then by exhaustion. On $E$, Lemma~\ref{lem:bv} gives $\ind_E\mathrm d\lambda=0$. Combining the restrictions proves \eqref{eq:lambdaODE} on the entire overlap.

Since $K$ is bounded away from zero, \eqref{eq:lambdaODE} has a bounded nonnegative density. Thus $\lambda$ is absolutely continuous, nondecreasing, and Lipschitz. The identities $t_s=b-M_s(\lambda)$ then imply absolute continuity of the inverse bids. Their exterior portions are affine, so the conclusion holds on their full supports. No assertion in the paper's Appendix A has been used to obtain this regularity.
\end{proof}

\begin{proposition}[Exact overlap length]
\label{prop:length}
Let $E=\{b\in[L,U]:\lambda(b)=1/2\}$. In a continuous, strictly increasing equilibrium with overlap, $E$ is an interval, possibly a single point, and
\begin{equation}
\label{eq:length}
 T=U-L=\Ts+\Leb(E).
\end{equation}
The lower and upper moving portions each have length $\Ts/2$.
\end{proposition}
\begin{proof}
Continuity, monotonicity, and the endpoint values $0,1$ imply that $E$ is a nonempty interval. By \eqref{eq:lambdaODE} and the chain rule,
\[
 \mathrm d(H(\lambda(b)))
 =K(\lambda(b))\,\mathrm d\lambda(b)
 =\ind_{E^c}(b)\dd b.
\]
Integration from $L$ to $U$ gives
\[
 H(1)-H(0)=(U-L)-\Leb(E).
\]
Appendix~\ref{app:algebra} evaluates $H(1)=\Ts$. Proposition~\ref{prop:endpoints} gives $U-L=T$, proving \eqref{eq:length}. The symmetry $K(1-\lambda)=K(\lambda)$ implies $H(1/2)=\Ts/2$, and integrating separately on the two moving portions proves the last assertion.
\end{proof}

\subsection{Relation to the original slope ratio}
The parameter cannot equal zero or one at an interior bid: by monotonicity this would create a nondegenerate constant interval at an endpoint level, contradicting \eqref{eq:lambdaODE}. Thus $0<\lambda<1$ in the interior. The derivative formulas below show that both inverse slopes are positive almost everywhere there. Writing
\[
 \rho=\frac{t_2'}{t_1'}=\frac{\lambda}{1-\lambda},
\]
the pivotal means become exactly the ones in the draft:
\[
 Q_1(\rho)=\frac{11\rho+9}{2(11\rho+13)},\qquad
 Q_2(\rho)=\frac{17\rho+11}{2(13\rho+11)}.
\]
The compact kernel satisfies
\[
 K(\lambda)=\frac{G(\lambda/(1-\lambda))}{(1-\lambda)^2},
\quad
 G(\rho)=\frac{22(169\rho^2+334\rho+169)}
 {(11\rho+13)^2(13\rho+11)^2}.
\]
The differential equation in the original notation is
\begin{equation}
\label{eq:rhoODE}
 (1-\rho)\,[G(\rho)\rho'-1]=0.
\end{equation}
The moving branch has $\mathrm db/\mathrm d\rho=G(\rho)$, but $\rho\equiv1$ on a bid interval is a distinct valid branch. Cancelling $1-\rho$ discards that branch. A removable singularity in a quotient is not a license to discard solutions of the original degenerate equation.

On moving portions, with $N(\lambda)=169-4\lambda+4\lambda^2$,
\begin{equation}
\label{eq:inverseslopes}
 t_1'(b)=\frac{48(1-\lambda)}{N(\lambda)},\qquad
 t_2'(b)=\frac{48\lambda}{N(\lambda)}.
\end{equation}
Both equal $1/7$ at $\lambda=1/2$ on a moving branch; on a genuine sticking interval they both equal one. This is compatible with finitely many nondifferentiable junctions in the direct bids.

\section{Construction and global verification of the continuous equilibria}
\label{sec:continuousconstruction}

\subsection{Construction for every \texorpdfstring{$T\ge\Ts$}{T at least Tstar}}
Put $\ell=T-\Ts\ge0$ and
\[
 a=\frac12+\frac{\Ts}{2},\qquad
 c=a+\ell=\frac12+T-\frac{\Ts}{2}.
\]
On the common-support interval $[1/2,T+1/2]$, construct inverse bids as follows:
\begin{align}
 b&=\frac12+H(\lambda),\quad
 t_s(b)=b-M_s(\lambda),
 &&0\le\lambda\le\tfrac12,\label{eq:constructlower}\\
 t_1(b)&=b-\frac5{12},\quad
 t_2(b)=b-\frac7{12},
 &&a\le b\le c,\label{eq:constructstick}\\
 b&=\frac12+\ell+H(\lambda),\quad
 t_s(b)=b-M_s(\lambda),
 &&\tfrac12\le\lambda\le1.\label{eq:constructupper}
\end{align}
For $\ell=0$, the middle interval is degenerate. Complete the direct bids outside the overlap by
\begin{align}
 b_1(t)&=t+\frac9{26},&&0\le t\le\frac2{13},\label{eq:lowerexterior}\\
 b_2(t)&=t+\frac{17}{26},&&T-\frac2{13}\le t\le T.\label{eq:upperexterior}
\end{align}
The values match at all joins by \eqref{eq:tasteendpoints} and $M_1(1/2)=5/12$, $M_2(1/2)=7/12$.

The moving inverse slopes in \eqref{eq:inverseslopes} are positive in the interior. On the sticking and exterior portions they equal one for the active signal. Thus each inverse is strictly increasing on its own support; a zero one-sided derivative at an endpoint does not create a constant interval. The direct strategies are continuous and strictly increasing, smooth on the finitely many interior pieces, and differentiable except at finitely many points. Their supports are
\begin{equation}
\label{eq:overlapsupports}
 I_1=[9/26,T+1/2],\qquad I_2=[1/2,T+17/26].
\end{equation}

\subsection{A global best-response criterion}
\begin{lemma}[Monotone surplus cutoffs]
\label{lem:globalsuff}
Suppose the opponent's bid distribution conditional on own signal $s$ is atomless, with finite positive measure $\nu_s$, and its quality-adjusted marginal surplus can be written as
\[
 t-\tau_s(b)
\]
for $\nu_s$-almost every $b$. If $\tau_s$ is nondecreasing on the occupied bid set and a prescribed bid $x_s(t)$ separates bids with $\tau_s(b)\le t$ from bids with $\tau_s(b)\ge t$, then $x_s(t)$ is a global best response. Intervals with zero bid probability may be crossed without changing payoff.
\end{lemma}
\begin{proof}
For $x>x_s(t)$, the payoff difference is the integral of $t-\tau_s(b)\le0$ over the additionally won bids. For $x<x_s(t)$, the payoff difference is minus the integral of $t-\tau_s(b)\ge0$ over the forgone wins. Neither deviation raises payoff. This argument includes deviations across support boundaries and across empty bid intervals.
\end{proof}

For the constructed continuous strategies, let
\[
 \tau_s(b)=b-\E[q\mid B_j=b,s_i=s]
\]
on the occupied bid set. On the overlap, the construction and its derivative ratios give $\tau_s=t_s$. Below the overlap only signal 1 can generate an opponent bid, and above it only signal 2 can do so. Therefore
\begin{equation}
\label{eq:globaltau}
\begin{array}{c|cc}
\text{bid range}&\tau_1(b)&\tau_2(b)\\ \hline
\bigl[9/26,1/2\bigr)&b-9/26&b-1/2\\
\bigl[1/2,T+1/2\bigr]&t_1(b)&t_2(b)\\
(T+1/2,T+17/26]&b-1/2&b-17/26
\end{array}
\end{equation}
These functions fit continuously and increase across the entire occupied bid range. Below signal 2's support, $\tau_2\le0$; above signal 1's support, $\tau_1\ge T$. Each prescribed bid crosses its own taste threshold. Lemma~\ref{lem:globalsuff} proves global optimality, not merely stationarity.

\subsection{Completion of the continuous classification}
Proposition~\ref{prop:separation} gives existence and necessity for separated supports. Propositions~\ref{prop:endpoints} and \ref{prop:length} show that overlap requires $T\ge\Ts$. The construction above provides equilibrium for every such $T$.

There is also uniqueness within the continuous class. In the separated case the pivotal equations force \eqref{eq:separatedbids}. In the overlapping case the endpoints are fixed, the lower moving branch is $b=1/2+H(\lambda)$, the upper moving branch is fixed by its upper endpoint, and \eqref{eq:length} fixes the intervening sticking length. Hence \eqref{eq:constructlower}--\eqref{eq:upperexterior} are forced.

For $T=1$, the sticking length is
\[
 1-\Ts=0.8205722574215637\ldots,
\]
and its endpoints are
\[
 a=0.5897138712892182\ldots,
 \qquad c=1.4102861287107818\ldots.
\]
This is why the original $T=1$ nonexistence claim fails. At $T=\Ts$ there is no sticking interval; this is not the only taste width admitting an exact equilibrium.

\section{An exact discontinuous equilibrium in the alleged nonexistence region}
\label{sec:jumps}

The previous results concern continuous direct bids. There is another escape from the differential path restriction besides atoms or singular inverses: \emph{both direct bids can jump across a common interval that no opponent ever bids in}. This section constructs such an equilibrium for every $2/13<T<\Ts$.

\subsection{Selecting the skipped mixture range}
The posterior means satisfy
\begin{equation}
\label{eq:reflection}
 M_1(1-\lambda)=1-M_2(\lambda),\qquad
 K(1-\lambda)=K(\lambda).
\end{equation}
For $0\le\alpha\le1/2$, set
\begin{align}
 g(\alpha)&=M_1(1-\alpha)-M_1(\alpha)
 =M_2(1-\alpha)-M_2(\alpha),\label{eq:g}\\
 \Gamma(\alpha)&=2H(\alpha)+g(\alpha).\label{eq:Gamma}
\end{align}
Here $g(\alpha)>0$ for $\alpha<1/2$, and
\begin{equation}
\label{eq:Gammabounds}
 \Gamma(0)=\frac2{13},\qquad \Gamma(1/2)=\Ts,
 \qquad
 \Gamma'(\alpha)=
 \frac{1056}{(13-2\alpha)^2(11+2\alpha)^2}>0.
\end{equation}
The derivative follows by differentiating \eqref{eq:Gamma} and using
$M_1'(1-\alpha)=M_2'(\alpha)$. Hence for each $T\in(2/13,\Ts)$ there is a unique $\alpha\in(0,1/2)$ with $\Gamma(\alpha)=T$.

\subsection{Generalized inverse cutoffs and direct bid jumps}
Fix this $\alpha$ and define
\begin{equation}
\label{eq:gapbounds}
 j_- =\frac12+H(\alpha),\qquad
 j_+=j_-+g(\alpha),\qquad
 \kappa_s=j_--M_s(\alpha).
\end{equation}
Equation~\eqref{eq:g} ensures
\begin{equation}
\label{eq:jumpmatch}
 j_+-M_s(1-\alpha)=\kappa_s,\qquad s=1,2.
\end{equation}
Define two continuous nondecreasing cutoff functions on $[1/2,T+1/2]$ by
\begin{align}
 b&=\frac12+H(\lambda),\quad
 t_s(b)=b-M_s(\lambda),
 &&0\le\lambda\le\alpha,\label{eq:jumplower}\\
 t_s(b)&=\kappa_s,
 &&j_-\le b\le j_+,\label{eq:jumpgap}\\
 b&=\frac12+T-\Ts+H(\lambda),\quad
 t_s(b)=b-M_s(\lambda),
 &&1-\alpha\le\lambda\le1.\label{eq:jumpupper}
\end{align}
At the start of the upper portion,
\[
 \frac12+T-\Ts+H(1-\alpha)
 =\frac12+T-H(\alpha)=j_+,
\]
using $T=2H(\alpha)+g(\alpha)$. Thus the pieces match exactly. The upper endpoint is $T+1/2$. Complete the exterior portions exactly as in \eqref{eq:lowerexterior}--\eqref{eq:upperexterior} and extend the cutoffs by zero and $T$ outside their ranges.

The cutoff functions are strictly increasing on the occupied pieces and constant on $[j_-,j_+]$. They are absolutely continuous. They correspond to direct strategies as follows: invert the lower occupied portions for $t<\kappa_s$ and the upper occupied portions for $t>\kappa_s$, and choose any bid in $[j_-,j_+]$ for the single taste $t=\kappa_s$. Each direct strategy is strictly increasing in taste, but has a jump from the left limit $j_-$ to the right limit $j_+$. The cutoff tastes lie strictly between zero and $T$. For example, $\kappa_2>0$ follows from the positive inverse slope on the initial moving portion; by reflection,
\begin{equation}
\label{eq:cutoffsum}
 \kappa_1+\kappa_2=T,
\end{equation}
so $\kappa_1<T$ as well.

No positive mass of types bids at the selected isolated cutoff bids. Thus the bid distributions are atomless, and $(j_-,j_+)$ is an empty bid interval, up to irrelevant zero-probability assignments.

\begin{proposition}[Discontinuous existence throughout the gap]
\label{prop:jump}
The strategies defined by \eqref{eq:jumplower}--\eqref{eq:jumpupper} and the exterior portions form a symmetric pure Bayesian Nash equilibrium for every $2/13<T<\Ts$ under the original second-price payment and symmetric tie-breaking rules.
\end{proposition}
\begin{proof}
On the two occupied overlap portions, the positive inverse slopes in \eqref{eq:inverseslopes} imply the mixture ratio $\lambda/(1-\lambda)$, so the actual conditional quality means are $M_s(\lambda)$. Therefore the quality-adjusted surplus cutoff is $\tau_s=t_s$ on each portion. On the exterior pieces, the first and last rows of \eqref{eq:globaltau} remain valid.

By \eqref{eq:jumpmatch}, the left and right limiting surplus cutoffs at the empty bid interval are the same number $\kappa_s$. Extend $\tau_s$ constantly through this interval. It is then nondecreasing throughout the total bid range. The interval itself has zero bid measure, so crossing it has no payoff effect. Each prescribed bid separates opponent bids with nonnegative marginal surplus from those with nonpositive marginal surplus. Lemma~\ref{lem:globalsuff} proves global optimality, including optimality of any gap bid for type $\kappa_s$. Ties have probability zero, so no modification of the mechanism is involved.
\end{proof}

\subsection{The explicit discontinuous equilibrium when \texorpdfstring{$T=1/6$}{T=1/6}}
The scalar equation $\Gamma(\alpha)=1/6$ has solution
\[
 \alpha=0.2497096895488589\ldots.
\]
The common skipped bid interval and the two cutoff tastes are
\begin{align*}
 j_-&=0.5450278788374292\ldots,&
 j_+&=0.6216387878292374\ldots,\\
 \kappa_1&=0.1650687526503747\ldots,&
 \kappa_2&=0.0015979140162920\ldots.
\end{align*}
These decimals are only illustrations; equations \eqref{eq:Gamma}--\eqref{eq:jumpupper} define the equilibrium exactly.

\begin{remark}[Why the earlier shorthand needs correction]
The statement ``no symmetric equilibrium with absolutely continuous inverse bids'' is too broad unless continuity and strict increase of the \emph{direct} bids remain part of the definition. The discontinuous equilibrium above has absolutely continuous \emph{generalized cutoff inverses}, constant on the skipped bid interval. What fails is the existence of continuous direct bids with connected individual bid ranges. Its escape route is neither pooling nor a singular inverse.
\end{remark}

\section{Audit of the previous draft and the remaining general question}
\label{sec:audit}

\subsection{What survives and what changes}
\paragraph{The exact calculations survive.}
The posterior means, the derivatives of $Q_s$, the factorization in \eqref{eq:rhoODE}, the integral $\Ts$, its symmetry, the antiderivative, and the moving inverse slopes are correct. The separated-support incentive threshold is $2/13$, not the geometric interval-disjointness threshold $4/13$. All these quantities have been derived explicitly above or in Appendix~\ref{app:algebra}.

\paragraph{The original $T=1$ counterexample fails.}
The missing $\rho=1$ sticking branch supplies length $1-\Ts$. Both bid functions remain strictly increasing there; it is the slope ratio, not either bid function, that is constant.

\paragraph{The revised $T=1/6$ counterexample is valid for the maintained class.}
It refutes Question~\ref{q:historical}, and Corollary~\ref{cor:elementary} does so without an absolute-continuity assumption. The full threshold classification follows from the measure argument and automatic regularity. Thus a singular inverse is not an unresolved loophole for this family.

\paragraph{The measure-theoretic patch can be simplified and strengthened.}
Working directly with $\rho\in[0,\infty]$ requires special care at zero, infinity, and bids of zero aggregate density. Proving that exceptional sets have Lebesgue measure zero alone does not automatically exclude every global continuation pathology. The compact parameter $\lambda$ and the signed-measure equation \eqref{eq:DAEmeasure} address all these issues at once, including genuinely singular inverse measures.

\paragraph{Unrestricted symmetric nonexistence does not follow.}
Proposition~\ref{prop:jump} supplies an exact discontinuous symmetric equilibrium for the same $T=1/6$ primitives. Thus the natural object of the negative result is \emph{continuous-strategy} existence. It is not necessary to change tie-breaking, introduce bid atoms, or invoke mixed strategies to obtain equilibrium in this particular example. Footnote 6 of \cite[p.~503]{PS2000} says that continuity was unnecessary for the authors' earlier binary-signal analysis; this should not be read as a theorem that every binary-signal equilibrium must itself have continuous bids.

\paragraph{The $\varepsilon$-equilibrium statement is an asymptotic one.}
Section IV of \cite[pp.~510--511]{PS2000} constructs continuous, strictly increasing approximate equilibria for sufficiently large $n$. It is not a theorem asserting such approximate equilibria for each fixed finite auction at every tolerance. The length mismatch here is not itself a proof or a complete explanation of their asymptotic construction.

\subsection{The abstract inverse-regularity issue}
Appendix A of the paper states that Assumption 2 implies the claimed differentiability and absolute continuity of inverse bids \cite[p.~513]{PS2000}. Such an implication is not valid for arbitrary strictly increasing smooth functions. For instance, let $C\subset[0,1]$ be a closed nowhere-dense set of positive Lebesgue measure, and set
\[
 f(x)=\int_0^x\operatorname{dist}(u,C)\dd u.
\]
Then $f$ is continuously differentiable and strictly increasing, with $f'=0$ on $C$. Since $f$ is Lipschitz and $f'=0$ on $C$, $f(C)$ has measure zero, while $f^{-1}(f(C))=C$ has positive measure. Thus $f^{-1}$ fails the null-set preservation property of absolutely continuous maps. This is an issue with a generic inference about inverses; it does not invalidate Proposition~\ref{prop:regularity}, which uses equilibrium optimality to prove regularity for the present family.

\subsection{Relation to later nonexistence work}
Jackson's nonexistence example \cite{Jackson2009} concerns common and private components, but uses finite private types, a binary common state, and signals some of which reveal the state exactly. It therefore is not a direct verification of nonexistence under the continuous-quality, positive-density, limited-information primitives used here. Its conclusion about symmetric equilibrium nonexistence is stronger in a different environment.

Araujo, de Castro, and Moreira \cite{ACM2004} study conditions for regular symmetric equilibria in their own class and, in another part of their analysis, use a modified second-price tie-breaking/payment procedure. Their mechanism-change result should not be read as an existence theorem for the unchanged mechanism in this note. Conversely, Proposition~\ref{prop:jump} does not change the mechanism at all.

\subsection{A precise broader question left by this analysis}
The singular-inverse question for the explicit family is answered here, and so is the existence of possibly discontinuous pure equilibria for every $T$ in that family. The remaining general existence issue is different.

\begin{question}[Symmetric existence without the continuity restriction]
\label{q:remaining}
For every primitive tuple $\mathcal A$ described in Section~\ref{sec:question}, is there a symmetric pure Bayesian Nash equilibrium under the original $(k+1)$st-price payment and symmetric tie-breaking rules, represented by measurable functions
\[
 b_s:[0,\bar t]\longrightarrow\R_+,
 \qquad t\le t'\ \Longrightarrow\ b_s(t)\le b_s(t'),
\]
where discontinuities and pooling intervals are allowed? In particular, does dropping Assumption 2 while retaining the paper's basic primitives always permit such an equilibrium?
\end{question}

The focus on pure, weakly taste-increasing strategies follows the reduction and best-response ordering discussed in footnotes 5 and 6 of \cite[p.~503]{PS2000}; it is not an assertion that such a reduction alone proves existence.

This is a continuation question not settled by the present proofs, not a claim that an exhaustive search has established its current literature status. The paper identifies the historical existence difficulty, and the related sources above explain why standard tie-breaking and multidimensional information matter. The explicit family here is \emph{not} a counterexample to Question~\ref{q:remaining}; it has pure symmetric equilibria for every $T$.

A complementary mathematical task is to characterize additional primitive conditions guaranteeing a \emph{continuous} equilibrium in the general model. For the uniform two-signal family, \eqref{eq:classification} is such a necessary-and-sufficient characterization.

\subsection{What the general differential-algebraic system says}
For reference, write the paper's factorization as
\[
 p(s\mid q,t)=\alpha_s(q)\beta_s(t).
\]
With clipped inverse tastes $t_s(b)$ and an absolutely continuous profile, define
\[
 Y(b\mid q)=\sum_r\alpha_r(q)\int_0^{t_r(b)}\beta_r(x)w(x)\dd x.
\]
The density of the $k$th-highest of the $n-1$ opponents' bids, conditional on $q$, is a positive constant times
\[
 Y^{n-k-1}(1-Y)^{k-1}\,\partial_bY.
\]
For a set of active signals $\mathcal R(b)$, at regular bids for which the corresponding own signal weights $\beta_s(t_s(b))$ are positive, pivotal optimality implies
\[
 A_{\mathcal R}(b,t)\,t'_{\mathcal R}(b)=0,
\]
where the active rows and columns have entries
\begin{align*}
 A_{sr}(b,t)={}&\beta_r(t_r)w(t_r)\\[-2pt]
 &\times\int_0^1[u(q,t_s)-b]f(q)\alpha_s(q)\alpha_r(q)
 Y^{n-k-1}(1-Y)^{k-1}\dd q.
\end{align*}
The factors involving the bidder's own taste and signal cancel from the posterior when they are positive. For a square active system with a nonzero derivative vector, singularity of $A_{\mathcal R}$ is necessary. One must use the correct active submatrix at a support face; inactive columns have zero derivative and cannot be treated as interior directions. A nonnegative null vector, consistent boundaries, and global incentive compatibility are additional requirements. None of these local conditions alone establishes a global continuous path. Empty bid intervals, as in Section~\ref{sec:jumps}, require a separate treatment and do not obey density-based pivotal equations in their interiors.

\appendix
\section{Proof of the bounded-variation facts}
\label{app:bv}

We provide the elementary one-dimensional facts used to remove the inverse-regularity assumption.

\subsection{Chain rule}
Let $v$ be continuous and of bounded variation on $[a,b]$, and let $F\in C^1$ near its compact range. The derivative $F'$ is uniformly continuous there. For a partition $a=x_0<\cdots<x_m=b$, Taylor's formula with a uniform modulus of continuity gives
\[
 F(v(x_i))-F(v(x_{i-1}))
 =F'(v(x_{i-1}))[v(x_i)-v(x_{i-1})]+R_i,
\]
with
\[
 \sum_i|R_i|
 \le\omega\!\left(\max_i|v(x_i)-v(x_{i-1})|\right)
 \Var(v;[a,b]),
\]
where $\omega(\varepsilon)\to0$. Uniform continuity of $v$ makes the right side vanish as the partition mesh tends to zero. The main sum converges to the Stieltjes integral of $F'(v)$ against $\mathrm dv$. Applying this on every subinterval gives
\[
 \mathrm d(F(v))=F'(v)\,\mathrm dv.
\]

\subsection{Derivative measures do not charge a level set}
Fix $c\in\R$. Choose continuous cutoffs $\chi_\varepsilon$ with $0\le\chi_\varepsilon\le1$, equal to one on $[c-\varepsilon,c+\varepsilon]$ and zero outside $[c-2\varepsilon,c+2\varepsilon]$. Define
\[
 F_\varepsilon(z)=\int_c^z\chi_\varepsilon(u)\dd u.
\]
Then $F_\varepsilon\in C^1$ and $\|F_\varepsilon\|_\infty\le2\varepsilon$. The chain rule gives
\[
 \mathrm d(F_\varepsilon(v))=\chi_\varepsilon(v)\,\mathrm dv.
\]
The left side tends to zero as a distribution because $F_\varepsilon(v)$ tends uniformly to zero. On the right side, dominated convergence with respect to the finite total variation measure $|\mathrm dv|$ gives convergence to $\ind_{\{v=c\}}\mathrm dv$, even in total variation. Hence that restricted signed measure is zero. Its total variation on the level set is therefore zero as well. This proves Lemma~\ref{lem:bv}.

\section{Algebra, integration, and endpoint checks}
\label{app:algebra}

\subsection{Conversion to the unbounded slope ratio}
With $\rho=\lambda/(1-\lambda)$, $M_s(\lambda)=Q_s(\rho)$. Differentiation gives
\[
 Q_1'(\rho)=\frac{22}{(11\rho+13)^2},\qquad
 Q_2'(\rho)=\frac{22}{(13\rho+11)^2}.
\]
The exact polynomial identity
\[
 (11\rho+13)^2-\rho(13\rho+11)^2
 =(1-\rho)(169\rho^2+334\rho+169)
\]
proves
\[
 Q_2'(\rho)-\rho Q_1'(\rho)=(1-\rho)G(\rho).
\]
The original slope-ratio equation follows by differentiating $t_s=b-Q_s(\rho)$ at regular bids and using $t_2'=\rho t_1'$.

\subsection{Integral and symmetry}
An antiderivative of $G$ is
\begin{equation}
\label{eq:antiderivative}
 F(\rho)=\frac{11}{288}\log\!\left(\frac{13\rho+11}{11\rho+13}\right)
 -\frac{11}{12(13\rho+11)}-\frac{13}{12(11\rho+13)}.
\end{equation}
Differentiating directly verifies $F'=G$. Its endpoint values give
\begin{align*}
 \int_0^\infty G(\rho)\dd\rho
 &=F(\infty)-F(0)\\
 &=\frac16+\frac{11}{144}\log\!\frac{13}{11}=\Ts.
\end{align*}
Since $G(1/\rho)/\rho^2=G(\rho)$,
\[
 \int_0^1G(\rho)\dd\rho
 =\int_1^\infty G(\rho)\dd\rho
 =\frac1{12}+\frac{11}{288}\log\!\frac{13}{11}.
\]
Equivalently, $K(1-\lambda)=K(\lambda)$ and $H(1/2)=\Ts/2$.

A compact closed form for $H$ that avoids the endpoint $\rho=\infty$ is
\begin{align}
 H(\lambda)={}&\frac{11}{288}
 \log\!\left(\frac{11+2\lambda}{13-2\lambda}\right)
 -\frac{11(1-\lambda)}{12(11+2\lambda)}
 -\frac{13(1-\lambda)}{12(13-2\lambda)}\nonumber\\
 &-\frac{11}{288}\log\!\frac{11}{13}+\frac16.
 \label{eq:Hclosed}
\end{align}
This is valid on all of $[0,1]$.

\subsection{Inverse slopes and the critical taste gap}
On a moving portion, $\lambda'=1/K(\lambda)$, so
\[
 t_s'=1-\frac{M_s'(\lambda)}{K(\lambda)}.
\]
Simplification gives \eqref{eq:inverseslopes}, or in the original parameter,
\[
 t_1'=\frac{48(\rho+1)}{169\rho^2+334\rho+169},\qquad
 t_2'=\frac{48\rho(\rho+1)}{169\rho^2+334\rho+169}.
\]
Also, with $z=4\lambda(1-\lambda)\in[0,1]$,
\[
 M_2-M_1=\frac{2(11+z)}{143+z},
\]
which is strictly increasing in $z$. This establishes the extrema $2/13$ and $1/6$ used in Corollary~\ref{cor:elementary}.

Finally,
\[
 \frac16-\frac2{13}=\frac1{78}>0,\qquad
 \Ts-\frac16=\frac{11}{144}\log\!\frac{13}{11}>0.
\]
Thus the proposed taste width lies strictly inside the continuous-equilibrium nonexistence interval.

\section{Reproducible checks and limits of the verification}
\label{app:checks}

The research audit recorded symbolic checks of the polynomial factorization, inverse slopes, antiderivative, reflection symmetry, kernel symmetry, and derivative of $\Gamma$. The recorded numerical diagnostics also constructed the separated, continuous-overlap, and discontinuous profiles and evaluated interim payoffs by quadrature with respect to their inverse-bid measures. The diagnostic code is not included with this manuscript.

Numerical best-response diagnostics were run at
\[
 T\in\{0.1,\ 2/13,\ 1/6,\ \Ts,\ 0.3,\ 1\}.
\]
For each signal, the program checks a grid of 31 tastes, adds the relevant jump-cutoff taste in the discontinuous case, and compares the prescribed payoff with 601 alternative bids. The moving-portion integrals use 64-point Gauss--Legendre quadrature; the affine portions are integrated analytically. No positive sampled deviation gain above $2\times10^{-10}$ was found.

These computations are diagnostics, not proofs over a continuum of types or bids. The mathematical justification is the signed-measure pivotal argument, the bounded-variation regularity proof, and the global sign-of-surplus verification in the text. The manuscript is not a proof-assistant formalization. The literature references establish the original model and related results, not priority for the calculations here.

\clearpage
\begin{thebibliography}{9}

\bibitem{PS2000}
Wolfgang Pesendorfer and Jeroen M. Swinkels (2000).
\newblock Efficiency and Information Aggregation in Auctions.
\newblock \emph{American Economic Review} \textbf{90}(3), 499--525.
\newblock DOI: \href{https://doi.org/10.1257/aer.90.3.499}{10.1257/aer.90.3.499}.
\newblock Source used: the published article. Relevant locations: existence discussion, pp.~501 and 505; assumptions, pp.~502--503; Proposition 1, p.~504; $\varepsilon$-equilibria, pp.~510--511; Appendix A, p.~513.

\bibitem{Jackson2009}
Matthew O. Jackson (2009).
\newblock Non-existence of equilibrium in Vickrey, second-price, and English auctions.
\newblock \emph{Review of Economic Design} \textbf{13}, 137--145.
\newblock DOI: \href{https://doi.org/10.1007/s10058-008-0059-2}{10.1007/s10058-008-0059-2}.
\newblock Author-hosted manuscript consulted:
\url{https://web.stanford.edu/~jacksonm/nonexist.pdf}.
\newblock The hosted version is dated 21 September 2005 and identifies itself as a revision of Jackson (1999). The journal citation above is to the published version.

\bibitem{ACM2004}
Aloisio Araujo, Luciano I. de Castro Filho, and Humberto Moreira (2004).
\newblock Pure Strategy Equilibria of Multidimensional and Non-Monotonic Auctions.
\newblock Working paper, March 2004, IMPA and EPGE/FGV.
\newblock \url{https://epge.fgv.br/files/1608.pdf}.
\newblock Relevant locations in the consulted manuscript: regular strategies, Definition 1, p.~11; modified second-price tie-breaking/payment procedure, pp.~22--23.

\end{thebibliography}
\end{document}
